\documentclass[10pt,a4paper]{amsart}
\usepackage[margin=19mm]{geometry}
\usepackage{amsmath,amssymb,mathtools}
\usepackage[scr=rsfs,cal=euler]{mathalfa}
\usepackage{xfrac}
\usepackage[unicode,hidelinks,bookmarks=false]{hyperref}
\newcommand{\ie}{i.e.\ }
\newcommand{\cf}{cf.\ }
\newcommand{\resp}{respectively\ }
\newcommand{\Subsection}{\subsection}
\providecommand{\nobreakdash}{\nobreak\hbox{-}}
\setlength{\emergencystretch}{2em}
\multlinegap0pt
\usepackage{bm}
\usepackage{bbm}
\usepackage{dsfont}
\usepackage{xspace}
\usepackage{cancel}
\usepackage{mathtools}
\usepackage{amsthm}
\makeatletter
\@ifundefined{definition}{\newtheorem{definition}{Definition}}{}
\@ifundefined{lemma}{\newtheorem{lemma}{Lemma}}{}
\@ifundefined{proposition}{\newtheorem{proposition}{Proposition}}{}
\makeatother
\makeatletter
\def\d{\mathrm{d}}
\newcommand{\dist}{d}%intrinsic distance
\newcommand{\dm}{d} % dimension
\newcommand{\kp}{\kappa}
\newcommand{\p}{o}% pole of M
\expandafter\ifx\csname remark\endcsname\relax
\newenvironment{remark}{{\flushleft {\bf Remark.}}}{}
\fi
\expandafter\ifx\csname remarks\endcsname\relax
\newenvironment{remarks}{{\flushleft {\bf Remarks.}}}{}
\fi
\makeatother
\title[Collective Optimization on Riemannian Manifolds with Bounded]{Collective Optimization on Riemannian Manifolds with Bounded Curvature}
\author{Hui Huang, Dohyun Kim,, Hansol Park}
\date{Source: arXiv:2606.14884v1 (CC BY 4.0)}
\begin{document}
\maketitle

\noindent\emph{Source and licence.} Collective Optimization on Riemannian Manifolds with Bounded Curvature. Version arXiv:2606.14884v1 (CC BY 4.0), distributed under the Creative Commons Attribution 4.0 International licence, as recorded in the arXiv licence metadata for this exact version and shown on the abstract page https://arxiv.org/abs/2606.14884v1. Full rights notice: licenses/arxiv-2606.14884-CC-BY-4.0.txt.\par\smallskip

\noindent\textit{Editorial scope.} Counted unit: the Proposition whose source label is \texttt{prop: positive mass} in the article (source0.tex lines 1351--1371, source label \texttt{prop: positive mass}), delivered together with the article's own complete original proof of it (source0.tex lines 1373--1590, one \texttt{proof} environment of 218 lines). No mathematics is written, repaired or re-derived: statement and proof are the article's own text line for line. Required context retained uncounted: Main-eq (lines 322--327); def:weaksol (lines 505--515); phi (lines 1284--1290); lem:lower (lines 1301--1307). Every definition, display and label of those lines is retained unchanged. Omitted from this package and not redistributed: 1--321, 328--504, 516--1283, 1291--1300, 1308--1350, 1372--1372, 1591--2465. Nothing inside a retained excerpt is omitted or altered: every retained body is delivered exactly as the article's lines are written, so no omission receipt of any kind accompanies this package. Citations: the retained excerpts carry no citation command at all, so no bibliography entry of the article is transcribed and no reference is redistributed. Reference adaptation: none was needed and none was made; every reference inside the delivered unit resolves inside it, so no unresolved reference or citation is produced. Printed slips kept literal: no printed word, symbol, hypothesis or number of the counted statement or of its proof was altered. Layout and class: the article's own class is \texttt{amsart} with packages including amsmath, amsfonts, amssymb, mathrsfs, graphicx, enumerate, mathabx, amscd, amsthm, bbm, mathtools, algorithm, algorithmic, subcaption; the wrapper is the lane's pinned \texttt{amsart} editorial wrapper at 10pt on A4, which restates the theorem-like environments the retained text needs and adds the article's own macro definitions that the retained text needs (\texttt{\textbackslash d}, \texttt{\textbackslash dist}, \texttt{\textbackslash dm}, \texttt{\textbackslash kp}, \texttt{\textbackslash p}), with extra wrapper packages (bm, bbm, dsfont, xspace, cancel, mathtools). The delivered numbering is the unit's own. Assets: no retained span contains a figure environment, a graphics-inclusion command, a PDF-page inclusion, a table or a TikZ picture; no image file is copied into this package, the package declares no asset dependency and no image file is redistributed with it. Licence: version arXiv:2606.14884v1 of this article is distributed under the Creative Commons Attribution 4.0 International licence, as recorded in the arXiv licence metadata for this exact version and shown on the abstract page https://arxiv.org/abs/2606.14884v1. A full rights notice is delivered at licenses/arxiv-2606.14884-CC-BY-4.0.txt. Characters: every retained excerpt is pure ASCII, so the delivered engine, XeLaTeX with its default Latin Modern text font, reports no missing character.
\begin{align}\label{Main-eq}
\begin{cases}
&\displaystyle\partial_t\rho_t(x)=-\lambda\mathrm{div}_M\left(u_\alpha^t(x)\rho_t(x)\right)+\frac{\sigma^2}{2}\Delta_M\left(h_{R-\delta, R}\left(\dist(x, \p)\right)^2\|u_\alpha^t(x)\|^2\rho_t(x)\right),\vspace{0.3cm}\\
&\displaystyle  u_\alpha^t(x)=\frac{h_{R, R+\delta}(\dist(x, \p))}{\int_{M}w_\alpha(y)\d \rho_t(y)}\int_{B_{R+\delta}(\p)}h_{R, R+\delta}(\dist(y, \p))w_\alpha(y)\log_xy~\d \rho_t(y),
\end{cases}
\end{align}

\begin{definition}[Weak solution] \label{def:weaksol}
Let $\rho_0\in\mathcal{P}(M)$ and $T>0$. We say $\rho\in C([0, T],\mathcal{P}(M))$ satisfies \eqref{Main-eq} with the initial data $\rho_0$ in the weak sense in the time interval $[0, T]$ if the following equality holds for all $\phi\in {C}_c^\infty(M)$ and all $t\in (0, T)$:
\begin{align*}
\frac{\d}{\d t}\int_M \phi(x)\d\rho_t(x)&=\lambda\int_{M} \nabla_M\phi(x)\cdot u_\alpha^t(x)\d\rho_t(x) \\
&\quad +\frac{\sigma^2}{2}\int_M h_{R-\delta,R}(\dist(x, \p))^2\|u_\alpha^t(x)\|^2\Delta_M\phi(x)\d\rho_t(x),
\end{align*}
and \[
\lim_{t\to 0+}\int_M \psi(x)\d\rho_t(x)=\int_M \psi(x)\d\rho_0(x)
\quad \text{for all } \psi\in C_b(M).
\]
\end{definition}

\begin{equation} \label{phi}
\phi_r(x) :=
\begin{cases}
\displaystyle e^{f(d(x,x^*))},\quad d(x,x^*)<r, \\
0,\quad \textup{otherwise}
\end{cases}, \quad f(s) := 1-\frac{r^2}{r^2 - s^2}.
\end{equation}

\begin{lemma} \label{lem:lower}
Let $\phi_r$ and $f$ be defined in \eqref{phi}. Then, we have
\begin{align*}
\Delta_M \phi_r(x)\geq\phi_r(x)\left(\frac{2r^2(3\dist(x, x^*)^4-r^4)}{(r^2-\dist(x, x^*)^2)^4}-\frac{2(\dm-1)r^2\sqrt{\kappa_-}\dist(x, x^*)}{(r^2-\dist(x, x^*)^2)^2}
\coth\left(\sqrt{\kappa_-}\dist(x, x^*)\right)\right).
\end{align*}
\end{lemma}

\begin{proposition}\label{prop: positive mass}
For all $t\in [0,T]$ and $r\in (0,R_0)$, we have
\begin{align*}
\int_{ x\in B_r(x^*)} \rho_t(x) \d x \geq  \left(\int_{x_\in B_r(x^*)} \rho_0(x) \d x\right) e^{-pt}
\end{align*}
for some $p>0$ given by
\[
p:= \max\left\{  p_1,p_2,\frac{2\lambda\beta}{m_r} \right\}.
\]
Here, $p_1$ and $p_2$ are defined in \eqref{p1p2}, $\beta$ is introduced in \eqref{K2} and satisfies  \eqref{const-beta}, and $m_r$ is defined in \eqref{mr}   for any $B>0$ with 
\[
\sup_{t\in [0,T]} \sup_{x\in B_R(\p)}\|u_\alpha^t(x) - \log_xx^*\|\leq B
\]
and for any $c\in (0,1)$ satisfying
\[
1+\sqrt{\kappa_-} r \leq \frac{3c^4-1}{2(d-1)(1-c)^2}.
\]

%\textcolor{blue}{HP: If we set $c_M=0$ then the condition on $c$ can be simplified.}

\end{proposition}

\begin{proof}
We recall
\[
\int_{x\in B_r(x^*)} \rho_t(x) \d x  \geq \int_M \phi_r(x) \rho_t(x) \d x .
\]
We use the formulation of a weak solution in Definition \ref{def:weaksol} to find
\begin{align*}
&\frac{\d}{\d t} \int_M \phi_r(x) \rho_t(x) \d x \\
&\quad = \int_M \phi_r(x) \partial_t \rho_t(x) \d x  \\
&\quad =\int_M   \phi_r(x) \left(   -\lambda \textup{div}_M (u_\alpha^t(x) \rho_t(x) ) + \frac{\sigma^2}{2} \Delta_M (  h_{R-\delta,R}(d(x,o))^2 \|u_\alpha^t(x)\|^2 \rho_t(x) )        \right) \d x \\ 
&\quad =   \int_M  \underbrace{  \lambda   \langle  \nabla \phi_r(x), u_\alpha^t(x) \rangle}_{=:T_1(x)} \rho_t(x) \d x        + \int_M \underbrace{ \frac{\sigma^2}{2}\Delta_M \phi_r(x) h_{R-\delta,R}\left(\dist(x, \p)\right)^2 \|u_\alpha^t(x)\|^2 }_{=: T_2(x)}\rho_t(x) \d x  \\
&\quad  = \int_M (T_1(x) + T_2(x) ) \rho_t(x) \d x .
\end{align*}
Our goal is to show that there exists $p>0$ such that 
\begin{equation} \label{T1T2}
T_1(x) + T_2(x) \geq -p\phi_r(x),\quad x\in M .
\end{equation}
Since the mollifier $\phi_r(x)$ vanishes outside the ball $\Omega_r := \{ x\in M : d(x,x^*) <r\}$, we restrict ourselves to $\Omega_r$.  From (A3)(i), we know that $B_{R_0}(x^*) \subseteq B_R(o)$. If $r\in (0,R_0)$, then
\[
d(x,x^*)<r <R_0 ~~ \Longrightarrow ~~x\in B_R(o) ~~ \Longrightarrow ~~d(x,o)<R.
\]
 We introduce two subsets of $\Omega_r$:
\begin{align*}
K_1 := \{ x\in M: d(x,x^*) >\sqrt{c}r\}
\end{align*}
and  
\begin{equation} \label{K2}
K_2 := \{ x\in M: - \lambda \langle \log_xx^*,u_\alpha^t(x)\rangle (r^2 - d(x,x^*)^2)^2 > \frac{r^2}{\beta}h_{R-\delta,R}^2 \|u_\alpha^t(x)\|^2 d(x,x^*)^2 \}
\end{equation}
where $\beta$ will be determined  later in \eqref{const-beta}. Then, we decompose $\Omega_r$ into three subsets
\[
\Omega_r = (K_1^c \cap \Omega_r) \cup (K_1\cap K_2^c \cap\Omega_r) \cup ( K_1\cap K_2\cap \Omega_r).
\]
Below,  we provide the desired estimate \eqref{T1T2} on each subset. \newline


\noindent $\bullet$ Case (i) (Subset $K_1^c \cap \Omega_r)$: In this subset, we have
\[
d(x,x^*)<\sqrt c r,\quad x\in K_1^c.
\]
For $T_1(x)$, we have
\begin{align*}
T_1(x)  & =  \lambda   \langle  \nabla \phi_r(x), u_\alpha^t(x) \rangle = \lambda \left\langle  \frac{2r^2 \log_xx^*}{ (r^2- d(x,x^*)^2)^2} \phi_r(x), u_\alpha^t(x)\right\rangle \\
& \geq -2r^2 \lambda \frac{  d(x,x^*) \|u_\alpha^t(x)\|      } {      (r^2- d(x,x^*)^2)^2   }  \phi_r(x)   \\
& \geq -\frac{ 2\lambda\sqrt c (B+\sqrt cr)}{(1-c)^2r} \phi_r(x) =: -p_1 \phi_r(x)
\end{align*}
where we used
\[
\|u_\alpha^t(x) \| \leq \| u_\alpha^t(x) - \log_xx^*\| + d(x,x^*) \leq B + \sqrt cr
\]
and
\[
\frac{1}{ r^2- d(x,x^*)^2  } <\frac{1}{(1-c)r^2}.
\]
On the other hand for $T_2(x)$,  we use the lower estimate of $\Delta_M \phi_r(x)$ in Lemma \ref{lem:lower} and the upper estimate  $h\leq1$ to find
\begin{align*}
&T_2(x)  = \frac{\sigma^2}{2}\Delta_M \phi_r(x)  h_{R-\delta,R}\left(\dist(x, \p)\right)^2 \|u_\alpha^t(x)\|^2 \\
& \geq \frac{\sigma^2}{2} \left(  \frac{2r^2(3d(x,x*)^4 - r^4)}{(r^2 - d(x,x*)^2)^4} - \frac{   2(d-1)r^2\sqrt{\kappa_-}d(x,x*)  }{(r^2 - d(x,x*)^2)^2}  \coth(\sqrt{\kappa_-} d(x,x*))\right) \phi_r(x) \|u_\alpha^t(x)\|^2 \\
&\geq -\frac{\sigma^2}{2}  \left(  \frac{2r^6}{ (r^2 - d(x,x^*)^2)^4} + \frac{2r^2 (d-1)}{ (r^2 - d(x,x^*)^2)^2} (1+\sqrt{\kappa_-}d(x,x^*)     \right) \phi_r(x) \cdot   \|u_\alpha^t(x)\|^2 \\
& \geq -\frac{\sigma^2}{2} \left(   \frac{2r^6}{(1-c)^4r^4}  + \frac{2r^2(d-1)}{(1-c)^2 r^2}(1+\sqrt{\kappa_-}\sqrt cr)     \right)   (B+\sqrt cr)^2 \phi_r(x) \\
& = : -p_2\phi_r(x).
\end{align*}
Hence in $K_1^c \cap \Omega_r$, we have
\[
T_1(x) + T_2(x) \geq -\max \{p_1,p_2\}\phi_r(x)
\]
where $p_1$ and $p_2$ are defined as
\begin{equation} \label{p1p2}
p_1= \frac{ 2\lambda\sqrt c (B+\sqrt cr)}{(1-c)^2r},\quad p_2 = \frac{\sigma^2}{2} \left(   \frac{2r^6}{(1-c)^4r^4}  + \frac{2r^2(d-1)}{(1-c)^2 r^2}(1+\sqrt{\kappa_-}\sqrt cr)     \right)   (B+\sqrt cr)^2 .
\end{equation}

\vspace{0.3cm}


\noindent $\bullet$ Case (ii) (Subset $K_1\cap K_2^c \cap \Omega_r$):  In this subset, we have
\[
d(x,x^*)>\sqrt cr 
\]
and  
\[
\lambda \langle \log_xx^*,u_\alpha^t(x)\rangle (r^2 - d(x,x^*)^2)^2 < \frac{r^2}{\beta} \|u_\alpha^t(x)\|^2 d(x,x^*)^2 .
\]
Our goal is now to show that 
\[
T_1(x) +T_2(x) \geq0, \quad x \in K_1\cap K_2^c \cap \Omega_r.
\]
 We observe
\begin{align*}
\frac{T_1(x) + T_2(x)}{2r^2\phi_r(x)}& =   \frac{  \lambda \langle \log_xx^*,u_\alpha^t(x)\rangle }{  (r^2- d(x,x^*)^2)^2 } + \frac{\sigma^2}{2} \Delta_M \phi_r(x) h_{R-\delta,R}\left(\dist(x, \p)\right)^2 \|u_\alpha^t(x)\|^2 \cdot \frac{1}{2r^2\phi_r(x)}.
\end{align*}
It follows from the lower estimate of $\Delta_{M}\phi_r(x)$ that 
\begin{align*}
 \frac{\Delta_M \phi_r(x)}{2r^2 \phi_r(x)} &\geq \frac{  3d(x,x^*)^4-r^4 }{   (r^2 - d(x,x^*)^2)^4 }    - \frac{  d-1  }{  (r^2-d(x,x^*)^2)^2} (1+\sqrt{\kappa_-} d(x,x^*)) \\
& = \frac{     (3d(x,x^*)^4-r^4  )  -(d-1)(1+\sqrt{\kappa_-}d(x,x^*)) (r^2 - d(x,x^*)^2)^2  }{    ((r^2 - d(x,x^*)^2)^4       } .
\end{align*}
%and from the upper estimate of $\eta$ that 
%\[
% \eta^2(\|u_\alpha^t(x)\|) \leq C_{2,\eta}^2 \|u_\alpha^t(x)\|^2.
%\]
Hence, we have
\begin{align*}
\frac{T_1(x) + T_2(x)}{2r^2\phi_r(x)} & \geq \frac{  \lambda \langle \log_xx^*,u_\alpha^t(x)\rangle (r^2- d(x,x^*)^2)^2 }{  (r^2- d(x,x^*)^2)^4 } \\
&\quad + \frac{\sigma^2  }{2} h_{R-\delta,R}\left(\dist(x, \p)\right)^2 \|u_\alpha^t(x)\|^2  \\
&\quad\quad \times \frac{     (3d(x,x^*)^4-r^4  )  -(d-1)(1+\sqrt{\kappa_-}d(x,x^*) (r^2 - d(x,x^*)^2)^2  }{  (  (r^2 - d(x,x^*)^2)^4     }  .
\end{align*}
In order to verify $T_1(x)+T_2(x)\geq0$, it suffices to show that
\begin{align*}
&\left(     \lambda \langle \log_xx^*,u_\alpha^t(x)\rangle   - \frac{\sigma^2  (d-1)}{2} h_{R-\delta,R}^2\|u_\alpha^t(x)\|^2 (1+\sqrt{\kappa_-} d(x,x^*) )\right) (r^2 - d(x,x^*)^2)^2\\
& \quad + \frac{\sigma^2   }{2}  h_{R-\delta,R}\left(\dist(x, \p)\right)^2\|u_\alpha^t(x)\|^2   (3d(x,x^*)^4-r^4  )  \geq 0
\end{align*}
which is equivalent to 
\begin{align*}
&\left(    -\lambda \langle \log_xx^*,u_\alpha^t(x)\rangle   + \frac{\sigma^2   (d-1)}{2} h_{R-\delta,R}^2\|u_\alpha^t(x)\|^2 (1+\sqrt{\kappa_-} d(x,x^*) )\right) (r^2 - d(x,x^*)^2)^2 \\
&\leq  \frac{\sigma^2    }{2}  h_{R-\delta,R}\left(\dist(x, \p)\right)^2\|u_\alpha^t(x)\|^2   (3d(x,x^*)^4-r^4  ).
\end{align*}
Our goal is to verify  that
\begin{equation} \label{case2-1}
 -\lambda \langle \log_xx^*,u_\alpha^t(x)\rangle(r^2 - d(x,x^*)^2)^2  \leq \frac{\sigma^2   }{4} h_{R-\delta,R}^2 \|u_\alpha^t(x)\|^2   (3d(x,x^*)^4-r^4  )
\end{equation}
and
\begin{align} \label{case2-2}
\begin{aligned}
  &\frac{\sigma^2   (d-1)}{2}h_{R-\delta,R}^2\|u_\alpha^t(x)\|^2 (1+\sqrt{\kappa_-} d(x,x^*) )(r^2 - d(x,x^*)^2)^2 \\
  &\quad \leq \frac{\sigma^2   }{4}h_{R-\delta,R}^2 \|u_\alpha^t(x)\|^2   (3d(x,x^*)^4-r^4  ).
\end{aligned}
\end{align}
 By adding these two inequalities \eqref{case2-1} and \eqref{case2-2}, we conclude that $T_1(x)+T_2(x)\geq0$. For the first inequality \eqref{case2-1}, we assume that $\beta$ introduced in subset $K_2$ satisfies
\begin{equation} \label{const-beta}
\frac{1}{c\beta} <\frac{\sigma^2  }{4}\left(3-\frac{1}{c^2}\right).
\end{equation}
Then, we have
\begin{align*}
 -\lambda \langle \log_xx^*,u_\alpha^t(x)\rangle(r^2 - d(x,x^*)^2)^2 &<  \frac{r^2}{\beta}h_{R-\delta,R}^2 \|u_\alpha^t(x)\|^2 d(x,x^*)^2 \\
 &< \frac{1}{c\beta}h_{R-\delta,R}^2 \|u_\alpha^t(x)\|^2 d(x,x^*)^4 \\
 &<\frac{\sigma^2  }{4}\left(3-\frac{1}{c^2}\right)h_{R-\delta,R}^2 \|u_\alpha^t(x)\|^2 d(x,x^*)^4  \\
 & < \frac{\sigma^2  }{4} h_{R-\delta,R}^2 \|u_\alpha^t(x)\|^2 (3d(x,x^*)^4 - r^4).
\end{align*}
Next, for the second inequality, it is equivalent to 
\[
(d-1)( 1+ \sqrt{\kappa_-}d(x,x^*)) (r^2 - d(x,x^*)^2)^2 \leq \frac{1}{2}  (3d(x,x^*)^4-r^4  ).
\]
Define an auxiliary polynomial
\[
g(s) := 3s^4 - r^4 - 2(d-1)(1+\sqrt{\kappa_-}s)(r^2- s^2)^2.
\]
Then, our goal is to show that $g(s)\geq0$. Since $\sqrt cr\leq d(x,x^*)=s\leq r$, we have
\begin{align*}
g(s) &\geq 3c^4 r^4-r^4 - 2(d-1)(1+\sqrt{\kappa_-}r)(1-c)^2 r^4 \\
&= ( 3c^4 -1 - 2(d-1)(1+\sqrt{\kappa_-}r))r^4 \geq0
\end{align*}
whenever
\[
1+\sqrt{\kappa_-} r \leq \frac{3c^4-1}{2(d-1)(1-c)^2}.
\]
Hence, for any $r>0$, we choose $c$ sufficiently close to 1 so that the above relation is satisfied. \newline


 

\noindent $\bullet$ Case (iii) (Subset $K_1 \cap K_2 \cap \Omega_r$): In this subset, we have
\[
d(x,x^*)>\sqrt cr
\]
and
\[
-\lambda \langle \log_xx^*,u_\alpha^t(x)\rangle (r^2 - d(x,x^*)^2)^2 > \frac{r^2}{\beta} h_{R-\delta,R}^2\|u_\alpha^t(x)\|^2 d(x,x^*)^2 .
\]
For $T_1(x)$, we recall the definition of $K_2$:
\begin{align*}
 \frac{   \langle \log_xx^*,u_\alpha^t(x)\rangle }{  (r^2- d(x,x^*)^2)^2 }  &\geq -\frac{ \| \log_xx^* \| \| u_\alpha^t(x) \|  }{  (r^2- d(x,x^*)^2)^2  } \geq  \frac{\lambda \beta}{r^2} \frac{  \langle \log_xx^*, u_\alpha^t(x)\rangle}{ \|\log_xx^*\| \|u_\alpha^t(x)\| h_{R-\delta,R}^2} \geq -\frac{\lambda \beta}{r^2 h_{R-\delta,R}^2}
\end{align*} 
Define
\[
D_r:= \sup_{x\in B_r(x^*)} d(x,o).
\]
Then, we claim that 
\[
D_r <R.
\]
Since $r<R_0$, we have $\overline{B_r(x^*)} \subseteq B_{R_0}(x^*) \subseteq B_R(o)$. Since $x\mapsto d(x,o)$ is continuous, it attains the maximum on $\overline{B_r(x^*)}$. Thus, there exists $x_{r_0} \in \overline{B_r(x^*)}$ such that
\[
d(x_{r_0},o) = \max_{x\in \overline{B_r(x^*)}} d(x,o).
\]
In addition, since $d(x_{r_0},o)<R$, we have
\[
D_r = \sup_{x\in \overline{B_r(x^*)}} d(x,o) \leq \max_{x\in \overline{B_r(x^*)}} d(x,o)  <R.
\]
This shows that since $d(x,o)\leq D_r<R$,  there exists $m_r>0$ such that 
\begin{equation} \label{mr}
h_{R-\delta,R}(d(x,o))^2 \geq m_r,\quad d(x,o)\leq D_r<R.
\end{equation}
Thus, we have
\begin{align*}
T_1(x)& =  \frac{ \lambda \langle \log_xx^*,u_\alpha^t(x)\rangle }{  (r^2- d(x,x^*)^2)^2 } 2r^2\phi_r(x)  \geq -2\frac{\lambda \beta}{m_r} \phi_r(x).
\end{align*}
Next, for $T_2(x)$, to see that $T_2(x)\geq0$, it suffices to show that
\[
 (3d(x,x^*)^4-r^4  )  -(d-1)(1+\sqrt{\kp_-}d(x,x^*)) (r^2 - d(x,x^*)^2)^2  \geq0 .
 \]
 In fact, in Case (ii), we show that
 \begin{align*}
 (3d(x,x^*)^4-r^4  ) &\geq 2(d-1)(1+\sqrt{\kappa_-}d(x,x^*) )(r^2 - d(x,x^*)^2)^2  \\
 &\geq (d-1)(1+\sqrt{\kappa_-}d(x,x^*)) (r^2 - d(x,x^*)^2)^2.
 \end{align*}
Thus, we get
\[
T_2(x)\geq0.
\]
Hence, we have
\[
T_1(x) + T_2(x) \geq -\max\left\{ p_1,p_2, \frac{2\lambda \beta}{m_r}\right\} \phi_r(x) =-p\phi_r(x)
\]
and this gives
\[
\frac{\d}{\d t} \int_M \phi_r(x) \rho_t(x) \d x  = \int_M (T_1(x) + T_2(x) ) \rho_t(x) \d x  \geq -p\int_M  \phi_r(x) \rho_t(x) \d x.
\]
Finally, Gr\"onwall's inequality gives the desired result. 
\end{proof}

\end{document}
